% Standalone problem document, generated from the adjacent README.md.
% Compile with XeLaTeX. Mathematical content is maintained in Markdown.
\documentclass[11pt,a4paper]{article}
\usepackage[fontsize=10.5pt]{scrextend}
\usepackage[margin=22mm,headheight=15pt,headsep=9mm,footskip=11mm]{geometry}
\usepackage{fontspec}
\setmainfont{texgyrepagella}[Extension=.otf,UprightFont=*-regular,BoldFont=*-bold,ItalicFont=*-italic,BoldItalicFont=*-bolditalic]
\setsansfont{texgyreheros}[Extension=.otf,UprightFont=*-regular,BoldFont=*-bold,ItalicFont=*-italic,BoldItalicFont=*-bolditalic]
\setmonofont{texgyrecursor}[Extension=.otf,UprightFont=*-regular,BoldFont=*-bold,ItalicFont=*-italic,BoldItalicFont=*-bolditalic,Scale=MatchLowercase]
\usepackage{amsmath,amssymb}
\usepackage{unicode-math}
\setmathfont{texgyrepagella-math.otf}
\usepackage{xcolor,microtype,enumitem,needspace,fancyhdr}
\usepackage{bookmark}
\definecolor{ink}{HTML}{17344B}
\definecolor{accent}{HTML}{197D80}
\definecolor{muted}{HTML}{596773}
\hypersetup{colorlinks=true,linkcolor=accent,urlcolor=accent,pdftitle={IS-03 - Johnson's
derivative-realizability
conjecture},pdfauthor={Open Problems in Numerical Linear Algebra}}
\urlstyle{same}
\setlength{\parindent}{0pt}
\setlength{\parskip}{5pt plus 1pt minus 1pt}
\linespread{1.04}
\setlength{\emergencystretch}{3em}
\widowpenalty=10000
\clubpenalty=10000
\displaywidowpenalty=10000
\allowdisplaybreaks[1]
\setlist{nosep,leftmargin=1.5em}
\providecommand{\tightlist}{\setlength{\itemsep}{0pt}\setlength{\parskip}{0pt}}
\providecommand{\pandocbounded}[1]{#1}
\pagestyle{fancy}
\fancyhf{}
\fancyhead[L]{\sffamily\footnotesize\color{muted}OPEN PROBLEMS IN NUMERICAL LINEAR ALGEBRA}
\fancyhead[R]{\sffamily\bfseries\footnotesize\color{accent}IS-03}
\fancyfoot[L]{\parbox[b]{0.9\textwidth}{\sffamily\fontsize{8}{10}\selectfont\color{muted}Editorial ratings; a bounded search cannot certify that a problem remains unsolved.\\Literature check: 2026-09-11}}
\fancyfoot[R]{\sffamily\footnotesize\color{muted}\thepage}
\renewcommand{\headrulewidth}{0.3pt}
\renewcommand{\headrule}{\hbox to\headwidth{\color{accent}\leaders\hrule height \headrulewidth\hfill}}
\makeatletter
\renewcommand{\section}{\Needspace{5\baselineskip}\@startsection{section}{1}{0pt}{12pt plus 2pt minus 2pt}{4pt}{\sffamily\bfseries\large\color{ink}}}
\renewcommand{\subsection}{\Needspace{4\baselineskip}\@startsection{subsection}{2}{0pt}{9pt plus 1pt minus 1pt}{3pt}{\sffamily\bfseries\normalsize\color{ink}}}
\makeatother
\setcounter{secnumdepth}{0}
\begin{document}
{\sffamily\small\color{muted}Eigenvalues and inverse problems\par}
\vspace{3pt}
{\raggedright\hyphenpenalty=10000\exhyphenpenalty=10000\sffamily\bfseries\fontsize{21}{24}\selectfont\color{ink}Johnson's
derivative-realizability conjecture\par}
\vspace{7pt}
{\sffamily\small\textbf{Difficulty:} challenging\quad\textbf{Importance:} interesting
to the community\par}
{\sffamily\small\bfseries\color{accent}Status: Solved\par}
\vspace{4pt}
{\color{accent}\hrule height 0.8pt}
\vspace{6pt}
\textbf{Rating rationale:} Challenging reflects preserving nonnegative
realizability under polynomial differentiation; community impact
connects the nonnegative inverse eigenvalue problem with polynomial
critical points.

\subsection{Resolution --- 2026-09-11}\label{resolution-2026-09-11}

\textbf{Negative resolution by Matthew J. Colbrook}, Department of
Applied Mathematics and Theoretical Physics, University of Cambridge.
\href{https://github.com/ajt60gaibb/OpenProblemsInNLA/blob/main/eigenvalues-and-inverse-problems/IS-03/solution.md}{Complete
manuscript} ·
\href{https://github.com/ajt60gaibb/OpenProblemsInNLA/blob/main/eigenvalues-and-inverse-problems/IS-03/solution.pdf}{PDF}
·
\href{https://github.com/ajt60gaibb/OpenProblemsInNLA/blob/main/eigenvalues-and-inverse-problems/IS-03/solution.tex}{LaTeX}.
\textbf{Theorem 1 and equations (1)--(7).}

The nonnegative real order-seven matrix
\(A=\operatorname{diag}(1/2,C_2,C_4)\) has a normalized
characteristic-polynomial derivative whose seventh power sum is
\(-8593/823543<0\). Every power of a nonnegative matrix has nonnegative
trace, so the derivative cannot be realized at order six, or after any
zero padding. Reducibility and positive trace are allowed in the
original target. This refutes its universal assertion.

The complete argument received an independent Codex-agent \textbf{PASS}
on 11 September 2026. The
\href{https://github.com/ajt60gaibb/OpenProblemsInNLA/blob/main/references/colbrook-additional-2026-09-11/verification/reviews/IS-03-review.md}{review
report} records the exact scope and a hash of the original reviewed
manuscript. The mathematical sections remain unchanged in the authored
version. Original ChatGPT generation is disclosed; no external human
peer review or formal proof certificate is asserted.
\href{https://github.com/ajt60gaibb/OpenProblemsInNLA/blob/main/references/colbrook-additional-2026-09-11/README.md}{Submission
and verification record}.

The original statement, source references and prior audit notes are
retained; the former difficulty rating is historical.

\subsection{Problem statement}\label{problem-statement}

For every integer \(n\geq5\) and every real entrywise-nonnegative matrix
\(A\in\mathbb R^{n\times n}\), define \(p_A(z)=\det(zI-A)\). Must there
exist an entrywise-nonnegative \(B\in\mathbb R^{(n-1)\times(n-1)}\) such
that

\[
\det(zI-B)=\frac1n p_A'(z)\qquad\text{as polynomials in }z?
\]

Equivalently, the critical points of the characteristic polynomial,
counted with multiplicity, would themselves form a realizable spectrum
of the smaller order. Neither \(A\) nor \(B\) is assumed symmetric or
diagonalizable. The realization must have exactly order \(n-1\);
allowing arbitrary additional zero eigenvalues changes the problem. This
would provide a dimension-reduction operation for nonnegative spectral
realization.

\newpage
\subsection{References}\label{references}

Hoover, McCormick, Paparella, and Thrall,
\href{https://arxiv.org/pdf/1712.05454}{\emph{On the realizability of
the critical points of a realizable list}}, Conjecture 1.2, p.~2, and
§6. The paper credits the conjecture to Johnson and records the
Cronin--Laffey low-order results.

\subsection{Earlier status check ---
2026-09-08}\label{earlier-status-check-2026-09-08}

Searches for
\textit{Johnson conjecture derivative nonnegative matrix characteristic polynomial proof counterexample},
\textit{1712.05454 2026}, and \textit{Monov conjecture solved} found
no general resolution. The source proves several classes and records the
solved cases \(n\leq4\), and \(n\leq6\) with \(\operatorname{tr}A=0\).
Nonnegative power sums alone are a different hypothesis. Monov's weaker
moment conjecture is not separately counted here.

\subsection{Audit update --- 2026-09-10}\label{audit-update-2026-09-10}

Rechecked the \href{https://arxiv.org/pdf/1712.05454}{primary
manuscript}, Conjecture 1.2 and its proved families; the
\href{https://doi.org/10.1016/j.laa.2018.06.024}{journal version} is LAA
555 (2018), 301--313. Results for Ciarlet/Suleĭmanova lists, appropriate
companion-matrix realizations, and trace-zero lists of orders five and
six settle substantive portions of the displayed target.
Johnson-conjecture/critical-point searches found no general proof or
counterexample.
\end{document}
